\documentclass{MMCStyle}
\graphicspath{{../figures/}}

\bianhao{202601139}
\tihao{A题}
\timu{长江十年禁渔下鱼类食物链恢复、混沌阈值与入侵风险的耦合建模}
\keyword{长江十年禁渔；四大家鱼；长江江豚；长江鲟；Hastings-Powell模型；李雅普诺夫指数}
\bibliographystyle{gbt7714-numerical}

\begin{document}
\begin{abstract}
围绕长江十年禁渔背景下鱼类生态系统的恢复与失衡问题，本文构建了“资源--功能群模块、珍稀物种模块和三级食物链混沌模块”三层框架。首先，用四类底层资源与四大家鱼功能群的耦合常微分方程刻画禁渔后资源与鱼群的传导；其次，用食源约束、人工放流和通道阻隔共同驱动的模型分析长江江豚与长江鲟的恢复过程；再次，用 Hastings-Powell 三级食物链模型展示稳定、周期振荡与更复杂不规则的长期行为，并在问题四中用最大李雅普诺夫指数识别混沌候选区间；最后，将工业污染和外来入侵物种纳入情景扩展。针对问题一，本文进一步引入四大家鱼总量与单网次捕获量代理（CPUE*）作为可观测指标，其中 2025 年“四大家鱼恢复到禁渔前约 1.8 倍”作为硬校准锚点，湖北段“单网次捕获量提升 120\%”仅作为拟合外的一致性检验；结果表明，单区校准后四大家鱼总量在 2025 年达到基线的 1.805 倍，CPUE* 达到 1.838 倍，对应相对误差分别为 0.26\% 和 16.47\%，而模型中的水草末值降至 0.204 倍，提示在当前参数组下恢复伴随底层承压上升。基线修复情景下，江豚和长江鲟末期规模分别达到初值的 1.133 倍和 1.140 倍；在污染与入侵叠加情景下，内部综合健康指数由 0.686 降至 0.575，且权重扰动后前两类情景排序存在交换，提示该指标并非完全稳健。据此，本文建议将“禁渔”与“水质治理、通道修复、定向增殖放流和入侵物种压制”联动实施。
\end{abstract}

\newpage
\tableofcontents
\newpage
\pagenumbering{arabic}
\pagestyle{fancy}

\section{问题重述}
长江流域是我国淡水渔业和生物多样性最集中的区域之一，但长期受到过度捕捞、工业废水排放、水域塑料污染以及葛洲坝、三峡大坝等工程的共同影响，鱼类洄游通道被切断，传统渔业资源和典型食物链持续退化。题面给出的背景信息表明，自 2021 年 1 月 1 日实施十年禁渔后，长江生态已经出现明显变化：局部水质改善，水草等水生植物恢复，青、草、鲢、鳙资源量在 2025 年监测中已恢复到禁渔前约 1.8 倍，说明禁渔对底层资源和经济鱼类具有显著修复作用。

但恢复并不意味着单向改善。题面同时指出，部分水域在禁渔后出现鱼类过快增长、过度摄食水草和浮游动物、导致水体浑浊和底栖植被减少的现象，甚至形成大面积“秃海”。与此同时，长江江豚数量回升至约 1240 头，中华鲟自然繁殖群体在 2024 年重新出现产卵场，外来入侵物种如鳄雀鳝、巴西龟又在鄱阳湖等支流形成稳定种群。由此可见，禁渔政策一方面推动了资源恢复，另一方面也可能诱发营养级失衡、通道恢复不足和入侵物种扩散等新问题，因此需要用统一模型区分“恢复”与“健康”的差别\cite{report2022,report2023}。

基于上述背景，本文围绕题面给出的五个子问题展开：第 1 问分析禁渔对水浮植物、浮游动物以及四大家鱼的影响；第 2 问研究禁渔收益是否能够传导到长江鲟类和长江江豚；第 3 问构建长江简单鱼类食物链并讨论长期演化；第 4 问判断三级食物链在不同承载力下是稳定、周期还是混沌；第 5 问则把工业废水、塑料污染和外来入侵物种纳入统一框架，评估十年禁渔对长江渔业生态的综合效应，并据此给出治理建议。题面中 2021--2026 年的不同年份信息，本文统一视为“政策起点--公报标定--现实验证”的连续口径来处理\cite{yangtzeA}。

\section{问题分析}
\subsection{总体分析}
五个问题不是孤立求值，而是沿着“底层资源--中层鱼群--珍稀物种--复杂食物链--综合情景”这条生态传导链展开。本文既要描述禁渔带来的恢复趋势，也要解释新失衡现象，并回答“为什么会恢复”“为什么会失衡”“何时会转向混沌”“污染和入侵会怎样放大脆弱性”。

\subsection{问题一分析}
问题一并不是简单比较禁渔前后的数量涨跌，而是要说明禁渔如何改变“底层资源--中层鱼类--上层捕食压力”之间的反馈方向。四大家鱼的食性不同，决定了它们对水草、浮游植物、浮游动物和底栖饵料的作用并不一致；因此，模型不能只给出末值比率，而应追踪恢复过程中的先后顺序、滞后关系和局地压力变化。这样才能解释同样处在禁渔背景下，为什么有些资源先回升，有些资源却会先被进一步挤压。

更具体地说，禁渔首先释放的是捕捞压力，随后才逐步改变种群之间的摄食关系。中层鱼类往往比底层资源更快恢复，这会带来更强的摄食和竞争效应，使局部水域出现“鱼类上升、底层下降”的阶段性现象。若只看末值，就会把“修复”和“失衡”混为一谈；只有把资源剩余、食压变化和恢复速度同时纳入，才能识别禁渔究竟是在推动系统回到更稳定的平衡，还是在形成新的局部承压。问题一的判据因此不应只是数量对比，还应包含能够反映资源压力和食压变化的综合指标。

从解释口径上看，这一问更适合按阶段来读：禁渔初期是捕捞压力迅速释放，鱼类数量先抬升；中期是中层鱼类对资源的摄食和空间竞争开始显现，底层资源的恢复被推迟；后期则要看系统能否回到新的平衡，而不是停留在局地高食压状态。也因此，问题一的结果应当用“资源压力是否回落、食压是否抬升、二者是否同步”来判断，而不是只写一个增长倍数。这样的分析既能解释底层资源为什么可能继续下滑，也能为后续问题二的珍稀物种响应提供上游输入。

因此，后文若出现底层资源暂时走低，也应结合时间滞后和局地压力一起解释，不能与修复结论直接对立。也只有把这一点写明，问题一的图和结论才不会只剩下单纯的数值对比，而是能形成“政策释放压力--食压重新分配--资源再平衡”这一完整链条。

\subsection{问题二分析}
问题二关心的是禁渔收益能否继续向珍稀物种传递，而不是单纯展示某个珍稀物种的数量是否增加。江豚和长江鲟都依赖食源恢复，但二者的约束条件并不相同：前者更接近对中层猎物供给的响应，后者则同时受到通道阻隔、栖息地破碎和人工放流的影响。若把它们都写成同一个增长率，模型就会把“自然恢复”和“补偿性增殖”混在一起，失去对机制差异的刻画。

因此，问题二更像是在搭建一个“珍稀物种响应模块”。题面附件中的事实需要逐一映射到参数：2022 年公报对应种群初值和监测基准，附件 1 中的航道阻隔、鱼道连通与栖息地破碎化对应 \(\delta_B\)，附件 2 中的放流、补充和食源恢复对应 \(R_S\)，而 \(e_D,m_D,H,u\) 则分别反映繁殖效率、自然死亡、捕捞/扰动和污染压力。只有把这些事实放进统一框架，才能说明禁渔收益是如何在食物网和生境连通性上被放大或削弱的，也才能解释为什么同处一个背景下，江豚和长江鲟会表现出不同的恢复曲线\cite{report2022}。

题面中的事实映射还意味着，问题二必须把“自然恢复”和“工程补偿”分开解释。对于江豚，更多是食源改善后的数量响应；对于长江鲟，除了食源，还要看洄游通道是否连续、放流是否提供额外补充、污染是否抵消恢复收益。因此，本问的重点不是比较哪个物种的绝对值更大，而是说明政策收益在不同约束条件下被保留了多少、损失了多少。把这层关系写清楚，才能让问题二从“珍稀物种是否增加”上升为“珍稀物种响应机制如何分化”。

表格的作用也在于让读者一眼看出每个参数到底来自哪条事实链，而不是把变量当成纯数学符号使用。这样写，问题二的证据链就会更完整：上游是禁渔和食源恢复，中间是通道、放流与污染修正，下游才是珍稀物种的数量变化。

\subsection{问题三分析}
问题三关注的是三层食物链在长期尺度上的行为类型，因此模型必须从复杂生态网中抽出一个仍然有代表性的核心结构。`Hastings-Powell` 三层模型并不是要精确复原长江流域的全部营养级，而是把“生产者--中层消费者--顶层捕食者”这一关键链条保留下来，再用承载力 \(K\) 把营养输入、水质改善和边界压力的综合效应压缩为一个可控尺度。这样，系统的长期行为就可以先被归纳为规则有界振荡、明显周期振荡或更复杂的动态区间。

这一问真正要回答的不是某个时刻的数值，而是“系统最后会落到哪一类状态”。当 \(K\) 较低时，链条更容易表现为规则有界振荡，说明资源供给不足以支撑更强的顶层反馈；当 \(K\) 增大时，捕食和增长之间的相互作用会被放大，系统可能转入明显的周期振荡，甚至进一步逼近更复杂的非线性行为。先把这些类型区分清楚，才能给后续问题四提供一个明确的参照：问题三回答的是“有没有复杂化”，问题四回答的是“复杂化是否已经跨入混沌”。

换言之，问题三并不要求把长江流域的真实结构全部装进模型，而是要找一个能够代表系统动力学方向的最小闭环。只要保住顶层捕食、底层供给和中层转移三者之间的非线性关系，就足以判断承载力变化会把系统推向哪类长期行为。这样处理的好处是，后面问题四的混沌判别就有了统一的参数轴，而不是在不同尺度上各说各话。对论文写作而言，这一问更适合以“行为分类”为主，不必把精力过度放在单个时点的曲线细节上。

同时，\(K\) 的变化并不是为了拟合某条曲线，而是为了识别承载能力变化对长期行为的方向性影响。因此，本问更关注区间上的动力学分型，而不是单点拟合精度；只要能稳定区分“稳定--周期--复杂”三类状态，就已经完成了问题三的任务。

\subsection{问题四分析}
问题四是在问题三的基础上进一步区分“看起来不规则”和“本质上混沌”。仅凭时间序列的起伏幅度，很难判断系统到底是长周期振荡还是混沌，因此需要引入最大李雅普诺夫指数作为判别量。若指数为负，扰动会被压回；若接近零，系统更接近周期轨道；若在扫描网格下持续为正，则说明初值差异会被指数放大，系统进入混沌可观测区间。

这里尤其要避免把某一个临界值写成绝对阈值。由于指数计算依赖步长、初值和网格密度，\(K\approx 3.22\) 更合适的表述应当是“在当前扫描分辨率下，正最大李雅普诺夫指数首次出现的起点区间”。因此，正文里既要保留粗扫结果，也要保留加密扫描结果，并把结论写成区间性的观测判断，而不是把一个网格点直接升格为硬边界。这样既尊重数值分析的不确定性，也能准确说明混沌是如何在参数轴上逐步显现出来的。

同时，参数扫描的写法也要服务于结论表达。粗扫负责定位区间，加密扫负责确认正指数并缩小范围；如果两者都指向同一段参数，就可以把这段写成“当前分辨率下的混沌起点区间”。这样的表述既避免把网格点误写成普遍阈值，也能提醒读者：这类数值结论依赖采样精度，属于当前模型和当前扫描设置下的观察结果，而不是脱离条件的绝对边界。也正因为如此，问题四的核心不是追求一个永远不变的临界数，而是给出一个足够稳健、可复核的混沌判定口径。

如果后续需要进一步细化，还可以通过更密的网格或不同初值重复验证，以说明结论对数值设置的稳健性。这样处理后，问题四的表述就会自然从“某个点的现象”转为“一个可复核的区间判断”，也更符合论文中对混沌识别的严谨写法。

如果把它放到图上展示，还应在曲线旁标明扫描范围和加密区间，避免读者把单个采样点误读成普适阈值。

\subsection{问题五分析}
问题五本质上是问题二和问题三的情景扩展：在禁渔基础上再叠加污染和外来入侵，考察多个扰动共同作用下的食物链响应。因此需要引入统一的综合评价口径，检验修复结论在污染和入侵持续存在时是否稳健。这里的综合健康指数只是内部比较指标，不是对外部生态监测的替代；它的作用，是把稀有物种、食源恢复和入侵压力放到同一标尺上，比较“仅禁渔”“禁渔加污染”“禁渔加污染加入侵”三种情景之间的相对位置，而不是追求一个绝对的生态真值。

由于污染会压低系统基础，外来入侵又可能通过竞争和掠食改变中层鱼类结构，问题五最需要说明的是指标是否稳健、排序是否稳定。若不同指标之间存在方向相反的变化，就必须通过权重敏感性说明来检验结论：在小幅权重扰动下，情景排序是否保持一致，综合指数是否仍能反映系统的真实差异。这样，问题五才能从单纯的数值汇总上升为对复合扰动机制的解释，即污染和入侵并不是简单叠加，而是通过削弱底层、挤压中层、改变顶层控制，进一步放大禁渔情景中的结构性差异。

再往下看，综合健康指数之所以只能是“内部指标”，是因为它需要服务于情景排序而不是替代生态监测。只要权重设计合理，并且在小幅扰动后排序不变，就说明结论并不是由某一个指标偶然拉动出来的；如果权重一变就翻盘，则说明模型只是把复杂现象压成了一个脆弱的分数，不能支撑治理判断。问题五真正需要交代的，是禁渔、污染和入侵三者如何共同作用于底层资源、中层鱼类和珍稀物种，并通过综合指数把这种联合作用转化为可比较的情景结果。

在治理建议层面，这也意味着污染控制和入侵防控不能分开看，而应与禁渔修复同步推进，才能保住综合收益。若只强调其中一项而忽略另外两项，综合指数的改善就可能停留在表面，难以转化为长期稳定的生态恢复。

\section{模型假设}
\begin{enumerate}
\item 在十年尺度内，长江流域宏观水环境可用归一化承载力表示，不对每个局地水域逐一建立空间网格。
\item 四大家鱼按照食性拆分为草鱼、鲢鱼、鳙鱼、青鱼四个功能群，分别对应水草、浮游植物、浮游动物和底栖饵料。
\item 江豚和长江鲟的恢复受食源供给约束，其中长江鲟额外依赖人工放流补给，且受大坝阻隔影响更强。
\item 污染通过降低承载力并提高死亡率起作用；外来入侵物种通过竞争或掠食中层鱼类改变营养级传导。
\item 为消除量纲不一致问题，本文所有状态变量均采用归一化尺度，初值为 1 附近时代表政策启动前后的相对基准量。
\end{enumerate}

\section{符号说明}
\begin{table}[htbp]
\centering
\caption{主要符号说明}
\begin{tabularx}{0.95\textwidth}{>{\raggedright\arraybackslash}p{2.2cm}X}
\toprule
符号 & 含义 \\
\midrule
$R_j(t)$ & 第 $j$ 类底层资源，分别表示水草、浮游植物、浮游动物、底栖饵料 \\
$C_j(t)$ & 第 $j$ 类四大家鱼功能群，分别表示草鱼、鲢鱼、鳙鱼、青鱼 \\
$M(t)$ & 中层饵料鱼或中层猎物供给 \\
$D(t)$ & 长江江豚种群规模 \\
$S(t)$ & 长江鲟种群规模 \\
$V(t)$ & 外来入侵物种强度 \\
$B_{\text{carp}}(t)$ & 四大家鱼总量，$G+S+B+Q$ \\
$R_{\text{basal}}(t)$ & 底层资源总量，$W+Ph+Z+N$ \\
$\mathrm{CPUE}^*(t)$ & 四大家鱼单网次捕获量代理指标（权重 1,1,1,0.7） \\
$H$ & 捕捞死亡率，禁渔场景下取 0，对照场景取正值 \\
$u$ & 污染强度参数 \\
$\delta_B$ & 通道阻隔和栖息地碎片化的综合损失 \\
$R_S$ & 人工放流强度 \\
$K$ & 三级食物链基础承载力参数 \\
$\Phi(t)$ & 食压指数 \\
$H_0$ & 反事实对照的初始捕捞压力参数（不参与观测校准） \\
$\lambda_{\max}$ & 最大李雅普诺夫指数 \\
$E(t)$ & 综合生态状态指数，由底层资源与四大家鱼共同构造 \\
\bottomrule
\end{tabularx}
\end{table}

\section{问题一的模型建立与求解}
\subsection{模型建立}
为避免把“湖北段局地观测”和“全流域平均口径”混为一谈，本文将问题一定位为流域平均单区模型，并用附件 2 的食性分工约束四大家鱼的功能群拆分。先给出观测锚点与变量映射表。
\begin{table}[htbp]
\centering
\caption{问题一的观测锚点与变量映射}
\label{tab:q1_anchor}
\begin{tabularx}{0.95\textwidth}{>{\raggedright\arraybackslash}p{2.2cm} >{\raggedright\arraybackslash}p{3.2cm} >{\raggedright\arraybackslash}p{3.2cm} >{\raggedright\arraybackslash}X}
\toprule
事实来源 & 观测口径 & 对应模型量 & 用途 \\
\midrule
题面与公报 & 2025 年四大家鱼恢复约 1.8 倍 & $B_{\text{carp}}(t)$ & 校准目标 \\
题面与公报 & 湖北段单网次捕获量提升 120\% & $\mathrm{CPUE}^*(t)$ & 局地一致性检验 \\
附件 2 & 草鱼、鲢鱼、鳙鱼、青鱼食性分工 & $C_1,\dots,C_4$ 与 $R_1,\dots,R_4$ & 机理约束 \\
附件 3 & 2022 年公报中的资源恢复事实 & 初值与时间锚点 & 规模标定 \\
\bottomrule
\end{tabularx}
\end{table}

在此基础上，本文采用资源--功能群配对模型：
\begin{equation}
\frac{\mathrm{d}R_j}{\mathrm{d}t}
=r_j R_j\left(1-\frac{R_j}{K_j}\right)-\alpha_j\frac{R_j C_j}{1+h_jR_j}-\eta_j u R_j,\quad j=1,\dots,4,
\end{equation}
\begin{equation}
\frac{\mathrm{d}C_j}{\mathrm{d}t}
=e_j\alpha_j\frac{R_j C_j}{1+h_jR_j}-(m_j+H)C_j.
\end{equation}
其中 $j=1,2,3,4$ 分别对应草鱼--水草、鲢鱼--浮游植物、鳙鱼--浮游动物、青鱼--底栖饵料。式中 $\eta_j u R_j$ 表示污染或环境扰动对底层资源的额外损失，$H$ 表示四类鱼共同受到的额外捕捞死亡率；在问题一中，禁渔情景取 $H=0$，对照情景取常态捕捞压力 $H=H_0>0$。因此，模型同时包含了“资源受压”和“鱼类受捕捞”的双向路径，能够刻画禁渔后资源层与消费者层的反馈关系，且本文结论针对当前 $H_0$ 设定成立。为量化局地承压风险，定义食压指数
\begin{table}[htbp]
\centering
\caption{问题一模型参数的角色与来源说明}
\label{tab:q1_param_role}
\small
\setlength{\tabcolsep}{4pt}
\begin{tabularx}{0.95\textwidth}{>{\raggedright\arraybackslash}p{2.8cm} >{\raggedright\arraybackslash}p{2.0cm} X >{\raggedright\arraybackslash}p{2.2cm}}
\toprule
参数 & 数值 & 来源/依据 & 角色 \\
\midrule
\multicolumn{4}{l}{\textbf{结构参数（归一化单区设定）}} \\
$r_w,r_p,r_z,r_n$ & \makecell[l]{$1.15,0.95,$\\$0.82,0.72$} & 反映水草、浮游植物、浮游动物和底栖饵料的相对恢复快慢次序；本文不将其解释为实测增长率，而是用于维持不同营养级的时间尺度分离。 & 时间尺度约束 \\
$K_w,K_p,K_z,K_n$ & $(1,1,1,1)$ & 统一固定为归一化承载上界，作为资源层的尺度基准，不解释为河段实测承载量。 & 归一化基准 \\
$\alpha_{wg},\alpha_{ps},\alpha_{zb},\alpha_{nq}$ & \makecell[l]{$1.22,0.92,$\\$0.84,0.76$} & 作为 Holling II 型摄食项的归一化耦合强度，取值保持摄食项与自增殖项同量级\cite{holling1959}。 & 结构耦合 \\
$h_w,h_p,h_z,h_n$ & \makecell[l]{$0.55,0.50,$\\$0.48,0.45$} & 用于控制饱和摄食的拐点位置，使系统在归一化尺度下保持有界响应。 & 饱和约束 \\
$e_{wg},e_{ps},e_{zb},e_{nq}$ & \makecell[l]{$0.43,0.35,$\\$0.33,0.30$} & 作为归一化转化效率，保持消费者摄食收益低于资源层自增殖量级，使四类鱼在禁渔与非禁渔情景下都维持合理数量级。 & 能量转化 \\
$m_g,m_s,m_b,m_q$ & \makecell[l]{$0.16,0.14,$\\$0.13,0.12$} & 作为归一化自然损失系数，与转化效率共同维持消费者层的背景损耗和平衡数量级。 & 背景损失 \\
\midrule
\multicolumn{4}{l}{\textbf{校准与情景参数（完整默认值见附录表 \ref{tab:q1_param_full}）}} \\
$consumer\_scale$ & $0.60$ & 固定缩放四大家鱼初值，保留既有物种占比，只承担单区归一化初值设定角色，不参与 2025 年观测拟合。 & 初值规范化 \\
$gain\_scale$ & $1.45$ & 对 2025 年四大家鱼总量锚点做单参数搜索，使 $B_{\text{carp}}(2025)\approx 1.8\,B_{\text{carp}}(2021)$。 & 唯一拟合参数 \\
$H_0$ & $0.25$ & 取自区间 $[0.10,0.40]$ 的中位值，代表中等捕捞压力，只用于构造“不禁渔”对照。 & 反事实对照 \\
$\mathrm{CPUE}^*$ 权重 & $(1,1,1,0.7)$ & 青鱼偏底栖、主要摄食螺类等底栖动物，在单网次指标中的可见性相对较低，因此作保守下调；该权重不参与单独估计。 & 一致性检验 \\
\bottomrule
\end{tabularx}
\normalsize
\end{table}
表 \ref{tab:q1_param_role} 仅保留正文所需的参数分层；问题一完整默认参数、场景控制量与初值模板统一列于附录表 \ref{tab:q1_param_full}，避免将关键数值仅保留在脚本中。
\begin{equation}
\Phi(t)=\frac{C_1(t)+C_2(t)+C_3(t)+C_4(t)}{R_1(t)+R_2(t)+R_3(t)+R_4(t)}.
\end{equation}
同时定义四大家鱼总量 $B_{\text{carp}}(t)=G(t)+S(t)+B(t)+Q(t)$、底层资源总量 $R_{\text{basal}}(t)=W(t)+Ph(t)+Z(t)+N(t)$，以及可观测的单网次捕获量代理
\begin{equation}
\mathrm{CPUE}^*(t)=G(t)+S(t)+B(t)+0.7Q(t),
\end{equation}
其中对青鱼的权重下调至 $0.7$，主要基于三点考虑：其一，青鱼偏底栖、主要摄食螺类等底栖动物，在单网次指标中的可见性通常低于草鱼、鲢鱼和鳙鱼；其二，本文的 $\mathrm{CPUE}^*$ 本质上是全流域平均口径下的观测代理，而不是逐网具标准化后的实测量，因此对青鱼采用保守折减比直接等权更稳妥；其三，该权重是经验假设，不参与单独估计。进一步地，在保持其余三类鱼权重为 $1$ 不变时，将青鱼权重在 $[0.5,1.0]$ 范围内扫描，2025 年 $\mathrm{CPUE}^*$ 倍率位于 $[1.805,1.862]$，不改变“湖北段事实只能作为局地一致性检验”的结论。由于分子与分母都由同类状态量构成，$\Phi(t)$ 仍仅作为内部压力指标；而 $\mathrm{CPUE}^*(t)$ 则用于对齐湖北段“单网次捕获量提升 120\%”这一局地一致性检验\cite{maunder2004,maunder2006}。换言之，问题一关注的不是单点末值，而是“观测锚点是否对齐、鱼群恢复是否快于底层恢复、资源是否在后期承压”这一完整链条。

\subsection{校准与锚点一致性检验}
数值上采用 \texttt{solve\_ivp} 积分，并将问题一的参数角色分成三层：一是维持归一化单区食物网结构的结构参数，二是仅服务于 2025 年总量锚点的校准参数，三是只用于“不禁渔”比较的反事实参数。具体做法是固定 $consumer\_scale=0.6$ 作为单区归一化初值设定，仅对增益缩放因子 $gain\_scale$ 做单参数搜索，使 2025 年四大家鱼总量对齐 1.8 倍这一硬锚点；$\mathrm{CPUE}^*$ 不参与目标函数，而作为拟合外的局地一致性检验。校准后得到 $gain\_scale=1.45$，此时 2025 年四大家鱼总量达到 1.805 倍，$\mathrm{CPUE}^*$ 达到 1.838 倍，对 1.8 倍与 2.2 倍的相对误差分别为 0.26\% 和 16.47\%。因此，本节对总量层面属于“硬锚点校准”，对湖北段 CPUE 事实则属于“校准后的外推一致性检验”，而非独立外部验证。随后取固定的中位反事实参数 $H_0=0.25$ 构造不禁渔情景；该参数位于预设区间 $[0.10,0.40]$ 的中部，不参与任何观测拟合，仅用于比较禁渔与持续捕捞的方向差异。在该设定下，2025 年四大家鱼总量和 $\mathrm{CPUE}^*$ 仍显著低于禁渔情景，且轨迹未出现反常上冲。

\subsection{结果与机制解释}
锚点一致性检验如表 \ref{tab:q1_result} 所示。与题面锚点相比，模型在总量层面已经较好贴合 2025 年恢复事实，而 CPUE* 代理虽低于 2.2 的局地软约束，但仍保持同向恢复，说明“湖北段 +120\%”更适合作为局地一致性检验而不是全流域平均的硬拟合目标。
\begin{table}[htbp]
\centering
\caption{问题一的 2025 年锚点一致性检验}
\label{tab:q1_result}
\begin{tabularx}{0.95\textwidth}{>{\raggedright\arraybackslash}p{3.0cm} >{\raggedright\arraybackslash}p{2.6cm} >{\raggedright\arraybackslash}p{2.6cm} >{\raggedright\arraybackslash}p{2.8cm}X}
\toprule
指标 & 观测值 & 模拟值 & 相对误差 & 说明 \\
\midrule
四大家鱼总量倍率 & 1.800 & 1.805 & 0.26\% & 对齐 \\
$\mathrm{CPUE}^*$ 倍率 & 2.200 & 1.838 & 16.47\% & 一致 \\
\bottomrule
\end{tabularx}
\end{table}

从过程上看，在当前参数组下，禁渔情景中四大家鱼总量与 $\mathrm{CPUE}^*$ 代理同步抬升，但模型中的底层资源末值仅为初值的 0.204 倍，水草最小值也出现在 10 年末，提示恢复并未同步改善底层资源。食压指数则从初值的 1 倍附近持续抬升至 4.213 倍，表明在该单区设定下鱼群恢复越强，底层资源承压越明显。结合附件 2 的食性分工，可将这一过程解释为：禁渔首先释放捕捞压力，随后中层鱼类对底层资源的摄食与竞争逐步增强，系统表现为“恢复与承压并行”而不是单向改善。

\begin{figure}[htbp]
\centering
\includegraphics[width=0.92\textwidth]{fig01_q1_foodweb.png}
\caption{问题一：校准后的禁渔与对照情景及观测锚点}
\label{fig:q1}
\end{figure}

\section{问题二的模型建立与求解}
\subsection{模型建立}
问题二不单独做外推拟合，而是把问题一得到的 $E(t)$ 作为上游食源代理，传入珍稀物种模块。这里 $M$ 表示中层猎物资源，$D$ 和 $S$ 分别表示江豚和长江鲟，$V$ 表示入侵竞争物种。为避免把自然恢复、通道修复和人工放流混在一起，先把题面事实与模型口径对应如表 \ref{tab:q2map} 所示。
\begin{table}[htbp]
\centering
\caption{问题二的观测锚点与模型映射}
\label{tab:q2map}
\begin{tabularx}{0.95\textwidth}{>{\raggedright\arraybackslash}p{2.4cm} >{\raggedright\arraybackslash}p{4.2cm} >{\raggedright\arraybackslash}p{3.6cm} X}
\toprule
事实来源 & 观测口径 & 对应模型量 & 角色 \\
\midrule
附件 1 & 洄游路径、大坝阻隔、鱼道连通性 & $\delta_B$ & 硬事实锚点 \\
附件 2 & 食性分工 & $a_D,a_S,a_V$ & 启发式先验 \\
附件 2 & 放流补偿 & $R_S$ & 情景量 \\
附件 3（2022 公报） & 江豚 1249 头、长江鲟 438 尾 & $D(0),S(0)$ & 硬事实锚点 \\
问题一输出 & 归一化生态指数 $E(t)$ & $E(t)$ & 内部代理 \\
\bottomrule
\end{tabularx}
\end{table}
其中只有 $\delta_B,D(0),S(0)$ 属于题面硬锚点；$a_D,a_S,a_V$ 是启发式先验，$R_S$ 是情景控制量，$E(t)$ 则是内部食源代理，不是外部观测值。
由问题一得到底层资源和四大家鱼轨迹后，构造综合生态状态指数
\begin{equation}
E(t)=0.62E_{\text{basal}}(t)+0.38E_{\text{carp}}(t),
\end{equation}
其中 $E_{\text{basal}}$ 和 $E_{\text{carp}}$ 分别由归一化的资源层和鱼群层加权平均得到。问题二将该指标作为上游食源输入，刻画江豚和长江鲟对禁渔收益的响应，而不是把珍稀物种部分写成独立的参数演示。进一步建立珍稀物种与入侵物种模型\cite{holling1959}：
\begin{equation}
\frac{\mathrm{d}M}{\mathrm{d}t}
=r_M M\left(1-\frac{M}{K_M(E)}\right)
-a_D\frac{MD}{1+h_MM}
-a_S\frac{MS}{1+h_MM}
-a_V\frac{MV}{1+h_MM}
-\lambda_u uM,
\end{equation}
\begin{equation}
\frac{\mathrm{d}D}{\mathrm{d}t}
=e_D a_D\frac{MD}{1+h_MM}-(m_D+\delta_B+\mu_Du)D,
\end{equation}
\begin{equation}
\frac{\mathrm{d}S}{\mathrm{d}t}
=R_S+e_S a_S\frac{MS}{1+h_MM}-(m_S+1.3\delta_B+\mu_Su)S,
\end{equation}
\begin{equation}
\frac{\mathrm{d}V}{\mathrm{d}t}
=e_V a_V\frac{MV}{1+h_MM}-(m_V+\mu_Vu)V.
\end{equation}
式中 $\delta_B$ 表示通道阻隔和栖息地碎片化的综合损失，$R_S$ 表示人工放流强度；题面事实与模型口径的对应关系已在表 \ref{tab:q2map} 中统一给出。

\subsection{情景结果与机制讨论}
以下结论仅在当前参数组与 2022 公报归一化口径下成立，不构成独立外部验证。在该口径下，江豚末值和长江鲟末值分别为 1.133 和 1.140。若仅提高通道阻隔而不加污染，江豚和长江鲟分别降至 0.901 和 0.856，显示附件 1 对应的通道约束本身就会压低恢复收益。进一步将长江鲟人工放流项置为 $R_S=0$ 时，长江鲟末值降至 0.519，而江豚仅变为 1.145，显示人工放流主要抬升鲟类，对江豚不是主解释项。污染情景中，江豚和长江鲟分别降至 0.623 和 0.706，提示污染会在高阻隔基础上继续侵蚀基线情景中的恢复收益。这里不把 $M$ 和 $V$ 放进结果表，是因为它们是传导和竞争状态，不是本问的直接观测目标。
\begin{table}[htbp]
\centering
\caption{问题二的结果口径、相对变化与对照说明}
\begin{tabularx}{0.95\textwidth}{>{\raggedright\arraybackslash}p{2.2cm} >{\raggedright\arraybackslash}p{3.6cm} >{\raggedright\arraybackslash}p{2.8cm} >{\raggedright\arraybackslash}p{1.7cm} X}
\toprule
项目 & 数值/情景输出 & 相对基线变化 & 是否外部验证 & 解释用途 \\
\midrule
2022 公报锚点 & 江豚 1249 头、长江鲟 438 尾 & -- & 是 & 作为问题二的初值与口径锚定 \\
基线情景末值 & 江豚 1.133，长江鲟 1.140 & -- & 否 & 作为禁渔与食源修复下的基线投影 \\
$R_S=0$ 对照 & 江豚 1.145，长江鲟 0.519 & 江豚 +1.0\%，长江鲟 -54.5\% & 否 & 拆分自然恢复与人工放流补偿效应 \\
高阻隔对照 & 江豚 0.901，长江鲟 0.856 & 江豚 -20.5\%，长江鲟 -24.9\% & 否 & 仅提高通道阻隔，分离生境约束效应 \\
污染情景末值 & 江豚 0.623，长江鲟 0.706 & 江豚 -45.0\%，长江鲟 -38.0\% & 否 & 比较污染叠加阻隔对基线收益的侵蚀 \\
\bottomrule
\end{tabularx}
\end{table}
这组对照表明，放流主要抬升鲟类，污染主要压低两类珍稀物种，因此问题二应收口为“情景比较与机制提示”，而不是独立外部验证。

\begin{figure}[htbp]
\centering
\includegraphics[width=0.88\textwidth]{fig02_q2_rare_species.png}
\caption{问题二和问题五共用的稀有种与入侵模块结果}
\label{fig:q2}
\end{figure}

\section{问题三的模型建立与求解}
\subsection{模型建立}
为判别长江简单食物链长期行为，采用 Hastings-Powell 三层模型\cite{hp1991}：
\begin{equation}
\frac{\mathrm{d}X_1}{\mathrm{d}t}
=rX_1\left(1-\frac{X_1}{K}\right)-\frac{a_1X_1X_2}{b_1+X_1},
\end{equation}
\begin{equation}
\frac{\mathrm{d}X_2}{\mathrm{d}t}
=\frac{a_1X_1X_2}{b_1+X_1}-\frac{a_2X_2X_3}{b_2+X_2}-d_1X_2,
\end{equation}
\begin{equation}
\frac{\mathrm{d}X_3}{\mathrm{d}t}
=\frac{a_2X_2X_3}{b_2+X_2}-d_2X_3.
\end{equation}
其中 $K$ 代表基础食物链有效承载力，可理解为营养输入、水质改善或污染冲击后的综合结果。

\subsection{求解与结果}
这里的 $K$ 是无量纲综合控制参数，不直接对应某一河段的实测承载力。固定当前参数组并取代表点后，$K\approx 1.8$ 附近对应规则有界振荡，$K\approx 2.7$ 附近对应明显周期振荡；更高 $K$ 下轨道会呈现更不规则的波动，但严格的混沌识别留到问题四通过最大李雅普诺夫指数完成。也就是说，问题三承担的是“代表性长期轨道展示”，而不是直接给出混沌判据。为避免只凭图像判读，这里统一在尾窗 $t\in[90,180]$ 上提取 $X_1$ 的峰谷差 $\Delta X_1$、平均峰间隔 $\bar T$ 和峰间隔变异系数 $\mathrm{CV}_T$ 作为观测型描述量，如表 \ref{tab:q3traj} 所示。

\begin{table}[htbp]
\centering
\caption{问题三代表轨道的尾窗描述量}
\label{tab:q3traj}
\small
\begin{tabularx}{0.95\textwidth}{>{\raggedright\arraybackslash}p{2.9cm} >{\centering\arraybackslash}p{1.2cm} >{\centering\arraybackslash}p{1.4cm} >{\centering\arraybackslash}p{1.4cm} >{\centering\arraybackslash}p{1.3cm} >{\centering\arraybackslash}p{1.5cm} >{\centering\arraybackslash}p{1.4cm}}
\toprule
类型 & $K$ & $\sigma(X_1)$ & $\Delta X_1$ & $N_{\text{peak}}$ & $\bar T$ & $\mathrm{CV}_T$ \\
\midrule
规则有界振荡 & 1.795 & 0.261 & 0.991 & 4 & 20.720 & 0.004 \\
明显周期振荡 & 2.665 & 0.502 & 2.092 & 3 & 34.178 & 0.014 \\
不规则样本 & 4.800 & 0.885 & 4.093 & 13 & 4.298 & 2.746 \\
\bottomrule
\end{tabularx}
\normalsize
\end{table}

表 \ref{tab:q3traj} 显示，随着 $K$ 增大，$X_1$ 的标准差由 0.261 升至 0.885，峰谷差由 0.991 扩大到 4.093，而峰间隔变异系数由 0.004 升至 2.746，说明系统从规则振荡逐步走向更不规则的长期行为。需要强调的是，第三行仅作为不规则样本展示，不在问题三中直接下混沌结论；是否进入混沌候选区间，仍以问题四的李雅普诺夫扫描为准。图 \ref{fig:q4} 的上半部分服务于问题三的代表轨道展示，下半部分的李雅普诺夫扫描属于问题四判据，不作为问题三单独的混沌证据。

\begin{figure}[htbp]
\centering
\includegraphics[width=0.96\textwidth]{fig03_q4_hastings_powell.png}
\caption{问题三共用代表轨道与问题四李雅普诺夫扫描}
\label{fig:q4}
\end{figure}

\section{问题四的模型建立与求解}
\subsection{模型建立}
通过对切向量方程做 Benettin 型迭代，计算最大李雅普诺夫指数 $\lambda_{\max}$。若 $\lambda_{\max}<0$，系统在当前积分窗口内更接近规则有界振荡；若 $\lambda_{\max}=0$ 附近，系统更接近周期轨道；若 $\lambda_{\max}>0$ 且在当前扫参网格下连续为正，则视为进入混沌候选区间。以下结果固定基准初值为 $(0.8,0.2,0.2)$，并采用步长 $dt=0.01$、瞬态截断 40、总积分时长 90 的数值设置。为了避免把局部数值现象误写成绝对阈值，文中再用相邻步长和微扰初值做一次括号复核。

\subsection{求解与结果}
在 1.4--4.8 的粗网格上计算最大李雅普诺夫指数，并在 3.0--3.4 之间加密扫描后，$\lambda_{\max}$ 首次连续转正出现在 3.2186 附近，并在 3.22--3.40 区间内保持为正。为检验这一结论在原始设定邻域内是否保留，进一步在 $K=3.20$ 与 $K=3.22$ 两侧做步长和初值微扰复核，如表 \ref{tab:q4robust} 所示；同时再做一轮更宽的压力测试，在 5 组初值、5 个步长和 3 个积分时长下共检查 75 组组合。

\begin{table}[htbp]
\centering
\caption{问题四近阈值复核}
\label{tab:q4robust}
\small
\begin{tabularx}{0.95\textwidth}{>{\raggedright\arraybackslash}p{4.1cm} >{\centering\arraybackslash}p{1.9cm} >{\centering\arraybackslash}p{1.9cm} >{\centering\arraybackslash}p{1.9cm} X}
\toprule
设置 & $\lambda_{\max}(3.20)$ & $\lambda_{\max}(3.22)$ & 判别 & 说明 \\
\midrule
基准步长 $dt=0.010$ & $-0.00135$ & $0.00117$ & 符号跨越 & 3.22 附近进入正指数区间 \\
较小步长 $dt=0.008$ & $-0.00134$ & $0.00117$ & 符号跨越 & 步长变化不改符号判断 \\
较大步长 $dt=0.012$ & $-0.00129$ & $0.00123$ & 符号跨越 & 步长变化不改符号判断 \\
微扰初值 $(0.8005,0.1995,0.2)$ & $-0.00087$ & $0.00165$ & 符号跨越 & 局部括号仍保留 \\
\bottomrule
\end{tabularx}
\normalsize
\end{table}

表 \ref{tab:q4robust} 说明：在基准初值、短积分窗 $t_{\text{total}}=90$ 和相邻步长扰动下，$K=3.20$ 到 $K=3.22$ 的符号跨越能够保留。但更宽的 75 组压力测试中，只有 5 组仍满足 $\lambda_{\max}(3.20)<0<\lambda_{\max}(3.22)$，且全部集中在基准初值与 $t_{\text{total}}=90$ 的情形；当积分时长延长到 150 或 200 时，对应的 50 组测试全部不再保留该符号跨越，其中多数转为两侧同时为正。需要说明的是，3.22 之前虽然出现过少量孤立正值，但它们未在加密网格下形成连续正区间，因此仅视为非稳健点，不作为混沌起点依据。故问题四更稳妥的结论应表述为“在固定初值和当前短窗求解设置下观察到的局部混沌候选括号约为 3.20--3.22”，而不是将 $K\approx 3.22$ 写成与求解设置无关的绝对硬阈值。图 \ref{fig:q4} 的上半部分仅供问题三共用轨道展示，下半部分与表 \ref{tab:q4robust} 一起构成问题四的局部判据证据。

\section{问题五的模型建立与求解}
\subsection{模型建立}
问题五直接复用问题二的状态方程，只在污染强度、入侵初值和竞争强度上做情景扩展；其中三组情景是内部构造的对照梯度，前两组共享相同的低入侵背景，第三组在污染基础上进一步提高入侵压力。这里继承的是状态方程和归一化口径，不继承外部验证口径。文中给出的综合健康指数是内部比较指标，用于三种场景的相对排序，不等同于外部监测口径；因此本问所有数值都应理解为内部情景输出，而不是外部生态健康验证。为了避免把单个终值误写成稳定性结论，本文将综合健康指数定义为尾窗平均：先对食源、江豚、长江鲟和入侵物种按各自基线归一化，再取末尾 $N_w=200$ 个离散步的加权平均。权重采用启发式内部赋权：食源、江豚、长江鲟和入侵物种的基准权重分别为 $0.52,0.28,0.20,-0.12$；敏感性分析中仅对食源权重做 $\pm 0.05$ 扰动，其余权重保持不变且不重归一化。记归一化后的中层鱼、江豚、长江鲟和入侵物种为 $\tilde M,\tilde D,\tilde S,\tilde V$，则内部健康指数可写为
\begin{equation}
H_{\text{int}}=\frac{1}{N_w}\sum_{k=T-N_w+1}^{T}\left(0.52\tilde M_k+0.28\tilde D_k+0.20\tilde S_k-0.12\tilde V_k\right),\quad N_w=200.
\end{equation}

\subsection{求解与结果}
三种情景“仅禁渔”“禁渔+污染”“禁渔+污染+入侵”的尾窗平均健康指数分别为 0.686、0.647 和 0.575，相比仅禁渔场景分别下降约 5.7\% 和 16.2\%。但这一排序并非完全稳健：当食源权重在 $\pm 0.05$ 范围内扰动时，前两类情景会发生交换，因此问题五更稳妥的结论是“污染和入侵都会压低系统健康度，且入侵叠加情景始终最差”，而不宜把“仅禁渔”和“禁渔+污染”的相对优劣写成稳定排序。
\begin{table}[htbp]
\centering
\caption{问题五的情景分解结果（末值与尾窗平均健康指数）}
\begin{tabularx}{0.95\textwidth}{>{\raggedright\arraybackslash}p{3.0cm} >{\centering\arraybackslash}p{1.6cm} >{\centering\arraybackslash}p{1.6cm} >{\centering\arraybackslash}p{1.6cm} >{\centering\arraybackslash}p{1.6cm} >{\centering\arraybackslash}p{1.6cm}}
\toprule
情景 & Food & Porpoise & Sturgeon & Invader & Health$_{\text{tail}}$ \\
\midrule
仅禁渔 & 0.380 & 1.133 & 1.140 & 0.568 & 0.686 \\
禁渔+污染 & 0.783 & 0.623 & 0.706 & 0.524 & 0.647 \\
禁渔+污染+入侵 & 0.673 & 0.608 & 0.701 & 0.625 & 0.575 \\
\bottomrule
\end{tabularx}
\end{table}
\begin{table}[htbp]
\centering
\caption{问题五的权重与敏感性说明（尾窗平均健康指数）}
\small
\setlength{\tabcolsep}{4pt}
\begin{tabularx}{0.95\textwidth}{>{\raggedright\arraybackslash}p{2.4cm} >{\raggedright\arraybackslash}p{3.8cm} >{\raggedright\arraybackslash}p{3.2cm} X}
\toprule
方案 & 权重 $(\tilde M,\tilde D,\tilde S,\tilde V)$ & 扰动口径 & 排序结果 \\
\midrule
基准权重 & $(0.52,0.28,0.20,-0.12)$ & 分量先归一化；权重不重归一 & 仅禁渔 \textgreater 禁渔+污染 \textgreater 禁渔+污染+入侵 \\
食源权重上调 & $(0.57,0.26,0.17,-0.12)$ & 仅调食源权重 $+0.05$ & 禁渔+污染 \textgreater 仅禁渔 \textgreater 禁渔+污染+入侵 \\
食源权重下调 & $(0.47,0.30,0.23,-0.12)$ & 仅调食源权重 $-0.05$ & 仅禁渔 \textgreater 禁渔+污染 \textgreater 禁渔+污染+入侵 \\
保守结论 & -- & 只保留跨权重稳定结论 & 入侵叠加情景始终最差；仅禁渔与禁渔+污染顺序不稳 \\
\bottomrule
\end{tabularx}
\normalsize
\end{table}
相对仅禁渔场景，污染情景的中层食物指标上升约 106.2\%，而江豚、长江鲟和综合健康指数分别下降约 45.0\%、38.0\% 和 5.7\%；在污染基础上再叠加入侵后，入侵物种指标继续上升约 19.3\%，综合健康指数再下降约 11.1\%。这说明问题五中最稳健的风险并不是“污染场景一定比仅禁渔更差多少”，而是“污染会削弱顶层恢复，而入侵叠加会把系统进一步推向最差情景”。同时，污染情景下中层食物指标由 0.380 升至 0.783，而综合健康指数却由 0.686 降至 0.647；这并不意味着生态更健康，而是与顶层受损一致的结构性偏移，提示中层猎物堆积和顶层控制减弱的可能性。

\begin{figure}[htbp]
\centering
\includegraphics[width=0.88\textwidth]{fig04_q5_scenarios.png}
\caption{问题五的情景比较结果}
\label{fig:q5}
\end{figure}

\section{模型的优缺点分析和优化}
\subsection{问题一模型的分析与改进}
本问的优点不在于继续增加模型维度，而在于把流域平均单区模型收紧成“结构参数设定--单硬锚点校准--局地一致性检验--反事实对照”的实证闭环。具体而言，2025 年四大家鱼恢复到禁渔前约 1.8 倍是唯一参与拟合的硬校准锚点；$consumer\_scale=0.6$ 作为归一化初值设定固定给定，不再与 2025 年锚点共同搜索；湖北段单网次捕获量提升 120\% 对应的 $\mathrm{CPUE}^*$ 仅作为拟合外的局地软约束。按这一口径，校准后四大家鱼总量倍率达到 1.805、相对误差为 0.26\%，而 $\mathrm{CPUE}^*$ 达到 1.838、相对误差为 16.47\%，说明模型在全流域总量层面已较好贴合题面事实，但对局地捕捞代理只能做到方向一致。与此同时，$H_0=0.25$ 仅作为区间 $[0.10,0.40]$ 内的固定中位反事实参数，不参与任何观测拟合；在该对照下，2025 年四大家鱼总量和 $\mathrm{CPUE}^*$ 分别为 0.716 和 0.731，且轨迹无反常上冲，因此“禁渔优于不禁渔”不再依赖目标导向的对照参数搜索。

$\Phi(t)$ 与水草末值的作用也应收紧理解：它们不是外部观测，而是用于解释“恢复与承压并行”的内部机制指标。当前参数组下，食压指数升至初值的 4.213 倍，水草末值降至 0.204 倍；对 \texttt{gain\_scale}、\texttt{consumer\_scale} 和 $H_0$ 做 $\pm10\%$ 扰动后，禁渔下鱼群恢复强于不禁渔、且系统承压上升这一主结论方向不变。其局限在于 $\mathrm{CPUE}^*$ 的 $(1,1,1,0.7)$ 权重属于经验假设，未单独估计，单区均质假设也无法直接解释湖北段等局地空间差异，因此本文不将湖北段观测写成全流域硬拟合结果。后续补强应优先增加独立观测锚点、统一代理口径，并继续补做更广范围的敏感性检验；空间异质性建模只能作为后续扩展，而不是当前版本的直接结论。

\subsection{问题二模型的分析与改进}
本问的优点不在于额外增加方程复杂度，而在于把题面事实、上游食源代理和政策对照统一到了同一套珍稀物种响应框架中。具体而言，2022 年公报中的江豚 1249 头和长江鲟 438 尾作为初值与口径锚点，附件 1 对应的通道阻隔事实通过 $\delta_B$ 进入模型，问题一输出的 $E(t)$ 仅作为内部食源代理输入，而 $a_D,a_S,a_V$ 与 $R_S$ 分别承担启发式先验和情景控制量的角色。正因为这种角色划分较清楚，问题二才能在基线、高阻隔、无放流和污染四组对照之间拆分出不同机制：基线情景下江豚和长江鲟末值分别达到 1.133 和 1.140；仅提高通道阻隔时二者降至 0.901 和 0.856，说明生境连通性本身就会压低恢复收益；将放流项置零后，长江鲟降至 0.519，而江豚仅约为 1.145，说明工程补偿主要抬升鲟类而非江豚。

其局限同样需要写清：问题二当前仍是 2022 公报归一化口径下的情景输出，而不是独立外部验证；其中 $E(t)$ 是内部代理，不应与外部监测量并列解读。另一方面，模型尚未引入年龄结构、季节性繁殖窗和分阶段放流调度，因此更适合用于比较“阻隔、放流、污染”三类因素的相对作用，而不宜外推为长期精细预测。后续若继续补强，应优先补入独立监测锚点和更细的放流/阻隔对照，再在此基础上引入年龄分层与繁殖季节结构，提高对自然恢复与工程补偿效应的拆分能力。

\subsection{问题三模型的分析与改进}
本问的优点不在于额外增加生态层级，而在于把规则有界振荡、明显周期振荡和不规则样本统一到了同一条无量纲控制参数轴上，并用尾窗描述量替代单纯的图像判读。具体而言，在当前参数组下，代表点 $K=1.795,2.665,4.800$ 分别对应规则有界振荡、明显周期振荡和不规则样本；与此同时，尾窗 $t\in[90,180]$ 上的 $\sigma(X_1)$ 由 0.261 升至 0.885，峰间隔变异系数 $\mathrm{CV}_T$ 由 0.004 升至 2.746，说明随着有效承载力提高，系统长期行为的波动幅度和不规则性都在增强。也正因为加入了这些可复核描述量，问题三的结论不再只是“看图判断”，而是形成了“参数点--代表轨道--尾窗特征”三者对应的展示链条。

其局限同样需要写清：这里的 $K$ 只是无量纲有效承载力控制参数，不直接对应某一河段的实测承载力，当前轨道分类仍建立在归一化参数组和有限时间窗之上，因此更适合作为代表性长期行为展示，而不是河段级定量识别。另一方面，第三行“不规则样本”只能说明轨道表现更复杂，并不在问题三中单独构成混沌判据；若只凭轨道图像，仍可能与长周期或有限窗数值现象混淆。因此，后续若继续补强，应优先引入实测营养级数据约束参数解释区间，并补充分岔图、相图和更多尾窗特征，而严格的混沌识别仍应继续由问题四的李雅普诺夫扫描承担。

\subsection{问题四模型的分析与改进}
本问的优点不在于单纯给出一个正的最大李雅普诺夫指数，而在于把问题三中“看起来更不规则”的轨道展示，进一步收敛为可复核的有限时间局部混沌候选括号。具体而言，在固定初值、步长和积分窗口的判据下，基准设置满足 $\lambda_{\max}(3.20)=-0.00135<0<0.00117=\lambda_{\max}(3.22)$；当基准初值保持不变、积分窗口固定为 $t_{\text{total}}=90$，并将步长扩展到 $0.005$、$0.008$、$0.010$、$0.012$、$0.015$ 时，这一符号跨越都能保留。因此，问题四确实能够把“看起来不规则”的轨道展示压缩为一个可复核的局部括号，但这个优点只在原始有限时间设置的邻域内成立。

其局限同样需要写清：这里的最大李雅普诺夫指数仍是固定初值、固定积分时长和固定扫参分辨率下的有限时间数值量，因此不能把当前结果外推为与数值设置无关的普适临界点。更关键的是，扩展压力测试在 75 组组合中只有 5 组保留 $\lambda_{\max}(3.20)<0<\lambda_{\max}(3.22)$，且一旦把积分时长延长到 150 或 200，符号跨越便不再出现，这说明该括号并不具备跨设置稳健性。另一方面，$3.22$ 之前虽然出现过少量孤立正值，但它们未形成连续正区间，也未被局部复核稳定支持，因此不能单独写成混沌起点。后续若继续补强，应优先在 $3.20--3.22$ 附近进一步加密扫参，并结合更长时间窗和分岔图做交叉验证；如果仍使用当前口径，则应把问题四明确理解为“基准有限时间设置下的局部候选括号识别”，而不是继续追求脱离条件的单点阈值。

\subsection{问题五模型的分析与改进}
本问的优点不在于再额外构造一套外部生态健康模型，而在于把问题二的归一化状态变量、污染压力和外来入侵强度统一到了同一套内部情景比较框架中，并用尾窗综合指标而非单点终值比较复合扰动效应。具体而言，综合健康指数 $H_{\mathrm{int}}$ 定义为加权归一化状态在尾窗 $N_w=200$ 上的均值，因此它承担的是内部排序口径而不是外部生态真值；三组情景也不是自然观测序列，而是构造的比较梯度：``仅禁渔''对应较低阻隔和零污染背景，``禁渔+污染''在相同入侵背景下提高阻隔并引入污染，``禁渔+污染+入侵''则进一步提高入侵参数并抬升初始入侵量。在基准权重下，三组健康指数分别为 0.686、0.647 和 0.575；相较于``仅禁渔''，污染使综合指数下降 5.72\%，而在污染基础上继续叠加入侵后又进一步下降 11.07\%，说明复合扰动会持续侵蚀禁渔收益。再结合对同一尾窗 $H_{\mathrm{int}}$ 做的权重敏感性检验，在仅对食源权重做 $\pm0.05$ 单因素扰动时，前两组情景会发生交换，但``禁渔+污染+入侵''始终最差，因此当前最稳妥的结论不是``三组排序完全稳定''，而是``禁渔+污染+入侵始终最差''这一保守判断。

其局限同样需要写清：首先，$H_{\mathrm{int}}$ 仍是基于问题二归一化状态变量的启发式内部赋权指标，不构成独立的外部生态健康验证，因此基准权重下``仅禁渔优于禁渔+污染''的结果只能视为当前口径下的比较输出，而不能写成普适结论。其次，当前敏感性分析只覆盖了同一尾窗指标上的单因素权重扰动，尚未系统展开更广的权重组合、污染强度和入侵强度不确定性，因此完整排序并不稳健，真正稳健的只有``加入侵情景最差''这一方向性结论。后续若继续补强，应优先引入独立生态健康监测指标并统一权重依据，再扩展到多权重、多污染水平和多入侵强度对照，从而把``最差情景稳健''与``完整排序稳健''明确区分开来。

\section{参考文献}
\bibliography{ref}

\appendix
\begin{center}
\heiti\zihao{4} 附\hspace{1pc}录
\end{center}
\section{问题一完整默认参数附表}
\begin{table}[htbp]
\centering
\caption{问题一完整默认参数附表}
\label{tab:q1_param_full}
\small
\setlength{\tabcolsep}{4pt}
\begin{tabularx}{0.97\textwidth}{>{\raggedright\arraybackslash}p{2.8cm} >{\raggedright\arraybackslash}p{2.2cm} X >{\raggedright\arraybackslash}p{2.2cm}}
\toprule
参数 & 默认值 & 依据/说明 & 角色 \\
\midrule
\multicolumn{4}{l}{\textbf{资源增长与归一化基准}} \\
$r_w,r_p,r_z,r_n$ & \makecell[l]{$1.15,0.95,$\\$0.82,0.72$} & 归一化相对恢复速度排序，表示水草快于浮游植物、浮游植物快于浮游动物、浮游动物快于底栖饵料；用于维持营养级时间尺度分离。 & 时间尺度约束 \\
$K_w,K_p,K_z,K_n$ & $(1,1,1,1)$ & 统一固定为归一化承载上界，不对应河段实测承载量。 & 归一化基准 \\
\midrule
\multicolumn{4}{l}{\textbf{摄食耦合与饱和响应}} \\
$\alpha_{wg},\alpha_{ps},\alpha_{zb},\alpha_{nq}$ & \makecell[l]{$1.22,0.92,$\\$0.84,0.76$} & Holling II 型摄食耦合强度，保持摄食项与资源自增殖项同量级\cite{holling1959}。 & 结构耦合 \\
$h_w,h_p,h_z,h_n$ & \makecell[l]{$0.55,0.50,$\\$0.48,0.45$} & 饱和摄食拐点参数，使系统在归一化尺度下保持有界响应。 & 饱和约束 \\
\midrule
\multicolumn{4}{l}{\textbf{转化效率与自然损失}} \\
$e_{wg},e_{ps},e_{zb},e_{nq}$ & \makecell[l]{$0.43,0.35,$\\$0.33,0.30$} & 归一化转化效率；沿“草鱼--水草、鲢鱼--浮游植物、鳙鱼--浮游动物、青鱼--底栖饵料”顺序设置为递减，避免消费者收益超过资源层可供给量级。 & 能量转化 \\
$m_g,m_s,m_b,m_q$ & \makecell[l]{$0.16,0.14,$\\$0.13,0.12$} & 归一化自然损失系数，与转化效率共同维持消费者层在禁渔和非禁渔情景下的数量级稳定。 & 背景损失 \\
\midrule
\multicolumn{4}{l}{\textbf{校准、反事实与场景控制}} \\
$consumer\_scale$ & $0.60$ & 固定缩放四大家鱼初值，保留既有物种占比，不由 2025 年观测拟合估计。 & 初值规范化 \\
$gain\_scale$ & $1.45$ & 由 2025 年“四大家鱼恢复到禁渔前约 1.8 倍”这一硬锚点单参数搜索得到。 & 唯一拟合参数 \\
$H_{\text{ban}}$ & $0$ & 禁渔基线情景下共同捕捞压力取零。 & 基线政策设定 \\
$H_0$ & $0.25$ & 不禁渔反事实情景的固定共同捕捞压力，取自区间 $[0.10,0.40]$ 的中位值，不参与观测校准。 & 反事实对照 \\
$u$ & $0$ & 问题一中污染强度默认取零，避免与问题五的污染情景混用。 & 基线环境设定 \\
\midrule
\multicolumn{4}{l}{\textbf{观测代理与初值模板}} \\
$\mathrm{CPUE}^*$ 权重 & $(1,1,1,0.7)$ & 草鱼、鲢鱼、鳙鱼等权；青鱼因偏底栖、单网次可见性较低而保守下调。 & 一致性检验代理 \\
$y_0^{\text{base}}=(W,Ph,Z,N,G,S,B,Q)$ & \makecell[l]{$(0.88,0.82,0.76,0.70,$\\$0.46,0.42,0.38,0.34)$} & 问题一的归一化初值模板；其中四类鱼初值在进入仿真前统一乘以 $consumer\_scale=0.60$。 & 状态初始化 \\
\bottomrule
\end{tabularx}
\normalsize
\end{table}
\section{主要源代码}
\lstinputlisting{../scripts/model_pipeline.py}

\end{document}
